\begin{lemma}
Let \(I\subseteq[n]\), put \(K=|I|\), and assume
\(\varepsilon_1\ge0\),
\(\noderef{ParameterHierarchyEventualComparisons}
(\varepsilon,\varepsilon_1,\varepsilon_2,n_0)\), \(n>n_0\),
\[
K=\noderef{TwoBiteNaturalIndependenceScale}(1+\varepsilon,n),
\]
and \(\noderef{FiberAndDegreeEvent}\).  Then the staged revelation touches at
most
\[
\varepsilon_1K^2
\]
projected pairs:
\[
|\noderef{TwoBiteTouchedProjectedPairs}(\omega,\varepsilon,I)|
\le \varepsilon_1K^2.
\]
\end{lemma}
\begin{proof}
Let
\[
P_I=\noderef{TwoBiteProjectionPairSet}(\omega,I),\qquad
T_I=\noderef{TwoBiteTouchedProjectedPairs}(\omega,\varepsilon,I).
\]
By the definition of \(\noderef{TwoBiteTouchedProjectedPairs}\), \(T_I\) is the
subcollection of \(P_I\) consisting of pairs satisfying
\(\noderef{TwoBiteProjectionPairTouched}\).  Unfold a member \(e\in P_I\).  If
\(e\) is red-tagged, then \(e=\{r,s\}\) with
\(r,s\in\pi_R(I)\); if \(e\) is blue-tagged, then
\(e=\{b,c\}\) with \(b,c\in\pi_B(I)\).  Thus every endpoint of a projected pair
in \(P_I\) satisfies \(\noderef{TwoBiteProjectionContains}(\omega,I,-)\).

Now unfold \(\noderef{TwoBiteProjectionPairTouched}\).  A touched red pair has
at least one red endpoint \(r\) for which
\(\noderef{TwoBiteStageRevealedVertex}(\omega,\varepsilon,I,\mathrm{inl}\,r)\)
holds, and similarly for blue pairs.  But
\(\noderef{TwoBiteStageRevealedVertex}\) is the disjunction
\[
\neg\noderef{TwoBiteProjectionContains}
\quad\text{or}\quad
\noderef{TwoBiteStageF1}
\quad\text{or}\quad
\noderef{TwoBiteStageF2}.
\]
The first alternative is impossible for an endpoint of a pair in \(P_I\), as
observed above.  Consequently every touched projected pair has at least one
endpoint in
\[
F_1\cup F_2,
\]
where \(F_1\) and \(F_2\) are the sets from
\(\noderef{TwoBiteStageF1}\) and \(\noderef{TwoBiteStageF2}\).

We first bound \(F_1\).  If \(x\in F_1\), then
\(\noderef{TwoBiteStageF1}\) gives both
\(\noderef{TwoBiteProjectionContains}(\omega,I,x)\) and
\[
|\,\noderef{TwoBiteBaseFiber}(\omega,x)\cap I\,|>\log n.
\]
The red fibers over \(\pi_R(I)\) are pairwise disjoint and partition \(I\) by
red projection; the blue fibers over \(\pi_B(I)\) are pairwise disjoint and
partition \(I\) by blue projection.  Therefore the total fiber mass over all red
projected vertices plus all blue projected vertices is exactly
\[
|I|+|I|=2K.
\]
Since every element of \(F_1\) is one of those projected vertices and contributes
more than \(\log n\) to this total mass,
\[
|F_1|\log n <2K,
\qquad\text{so}\qquad
|F_1|\le \frac{2K}{\log n}.
\]

Next consider \(F_2\).  If \(x\in F_2\), then
\(\noderef{TwoBiteStageF2}\) gives
\(\noderef{TwoBiteProjectionContains}(\omega,I,x)\) and says that the set of
same-color base neighbors \(y\) satisfying \(y\in F_1\) has cardinality more
than \(t_2/\log n\), where
\[
t_2=\noderef{TwoBiteLargeCutoff}(\varepsilon,n)=n^{1/4+\varepsilon}.
\]
Indeed, \(\noderef{TwoBiteBaseAdj}\) is false across colors, so these neighbors
have the same color as \(x\).  For each such neighbor \(y\in F_1\), the
\(\noderef{TwoBiteStageF1}\) condition gives
\[
|\,\noderef{TwoBiteBaseFiber}(\omega,y)\cap I\,|>\log n.
\]
As \(y\) varies over same-color base vertices, these fibers are disjoint.  They
all lie in \(\noderef{TwoBiteX}(\omega,I,x)\): if \(x\) and \(y\) are red, then
every vertex in the red fiber over \(y\) has red projection \(y\), and the
base-edge \(x\sim y\) puts it in the lifted red neighborhood of \(x\); the blue
case is identical.  Therefore
\[
|X_x(I)|> \frac{t_2}{\log n}\cdot \log n=t_2.
\]
Since \(X_x(I)\subseteq I\), either
\(|X_x(I)|\le\noderef{TwoBiteHugeCutoff}(n)\) or
\(\noderef{TwoBiteHugeCutoff}(n)<|X_x(I)|\).  In the first case the preceding
strict lower bound \(t_2<|X_x(I)|\) puts \(x\) in
\(\noderef{TwoBiteLargeClass}(\omega,\varepsilon,I)\); in the second case the
upper bound \(X_x(I)\subseteq I\) puts \(x\) in
\(\noderef{TwoBiteHugeClass}(\omega,I)\).  Thus
\[
x\in L_I\cup H_I=\noderef{TwoBiteLargeHugeVertices}(\omega,\varepsilon,I).
\]
Hence
\[
|F_2|\le |\noderef{TwoBiteLargeHugeVertices}(\omega,\varepsilon,I)|.
\]
The child lemma \(\noderef{LargeHugeVertexCountBound}\) has exactly the same
rounded scale hypothesis \(K=\noderef{TwoBiteNaturalIndependenceScale}
(1+\varepsilon,n)\), the same \(n>n_0\), the same eventual-comparisons package,
and the same \(\noderef{FiberAndDegreeEvent}\).  Applying it gives
\[
|F_2|<n^{1/4}.
\]

A single projected base vertex can be an endpoint of at most \(K\) pairs in
\(P_I\).  If the endpoint is red, its possible partners are contained in
\(\pi_R(I)\), whose cardinality is at most \(|I|=K\); if it is blue, its possible
partners are contained in \(\pi_B(I)\), again of cardinality at most \(K\).
Charge each touched pair to one of its endpoints in \(F_1\cup F_2\).  This may
double-count pairs, but it gives the upper bound
\[
|\noderef{TwoBiteTouchedProjectedPairs}(\omega,\varepsilon,I)|
\le K(|F_1|+|F_2|)
\le K\left(\frac{2K}{\log n}+n^{1/4}\right).
\]
Finally, the \(n>n_0\) instance of
\(\noderef{ParameterHierarchyEventualComparisons}\) is stated with
\[
k=(\noderef{TwoBiteNaturalIndependenceScale}(1+\varepsilon,n):\mathbb R)=K.
\]
One of its conjuncts is precisely
\[
K\left(\frac{2K}{\log n}+n^{1/4}\right)\le\varepsilon_1K^2.
\]
Combining this numerical inequality with the preceding display proves
\[
|\noderef{TwoBiteTouchedProjectedPairs}(\omega,\varepsilon,I)|
\le \varepsilon_1K^2,
\]
which is the asserted bound.
\end{proof}
